% Direct-authority diplomatic French for printed p.165.
\SourcePage{165}

\LayerTarget{tohoku:remark:3.4.2:1:p2}
donnée dans la démonstration de prop. 3.4.2, pour qu'on retrouve l'opérateur
bord usuel de la suite exacte de cohomologie relative à la suite exacte
$0\longrightarrow F\longrightarrow G\longrightarrow H\longrightarrow0$, dans
le cas où $G$ (donc aussi $F$ et $H$) est abélien. Dans ce cas, on a un
diagramme de suites exactes horizontales et verticales

\LayerTarget{tohoku:diagram:3.4.2:exact-sequence-grid}
\[
\begin{array}{ccccccccc}
&&0&&0&&0&&\\
&&\big\downarrow&&\big\downarrow&&\big\downarrow&&\\[0.45em]
0&\longrightarrow&F&\longrightarrow&G&\longrightarrow&H&\longrightarrow&0\\
&&\big\downarrow&&\big\downarrow&&\big\downarrow&&\\[0.45em]
0&\longrightarrow&\overline F&\longrightarrow&\overline G&\longrightarrow&\overline H&\longrightarrow&0\\
&&\big\downarrow&&\big\downarrow&&\big\downarrow&&\\[0.45em]
0&\longrightarrow&F'&\longrightarrow&G'&\longrightarrow&H'&\longrightarrow&0\\
&&\big\downarrow&&\big\downarrow&&\big\downarrow&&\\[0.45em]
&&0&&0&&0&&
\end{array}
\]
d'où on conclut [6, prop. III 4.1] que le diagramme d'homomorphismes cobords
suivant est \textit{anticommutatif} :

\LayerTarget{tohoku:diagram:3.4.2:boundary-anticommutative-square}
\[
\begin{array}{ccc}
H^0(X,H')&\longrightarrow&H^1(X,F')\\
\big\downarrow&&\big\downarrow\\[0.45em]
H^1(X,H)&\longrightarrow&H^2(X,F).
\end{array}
\]

\LayerTarget{tohoku:eq:3.4.2:dQ-minus-boundary-cE}
Le premier homomorphisme vertical est surjectif, et on vérifie que dans la
construction faite plus haut, $d(Q)$ s'obtient à partir de la classe
$c(E)\in H^1(X,H)$ par passage au quotient dans l'homomorphisme composé
\[
H^0(X,H')\longrightarrow H^1(X,F')\longrightarrow H^2(X,F),
\]
et est donc égal à $-\partial c(E)$.

\LayerTarget{tohoku:remark:3.4.2:2:local-triviality}
2. On n'oubliera pas que pour appliquer prop. 3.3.1, à un fibré algébrique
projectif, il faut d'abord vérifier qu'il est localement trivial.
Malheureusement, cette vérification, quand on ne sait pas à priori qu'on peut
remonter le groupe structural semble souvent difficile.

\LayerTarget{tohoku:example:projective-bundle-vector-lift}
Exemple : Sur une variété projective complexe sans singularités, on a un fibré
\textit{holomorphe} $E$ à fibre l'espace projectif (on sait, d'après
Kodaira--Borel, que $E$ est aussi une variété algébrique), provient-il d'un
fibré holomorphe vectoriel? La réponse est affirmative d'après prop. 3.4.1. si
$E$ est localement trivial du point de vue algébrique, i. e. si par tout point
de la base passe une section rationelle ; la réciproque est d'ailleurs vraie,
car [16] tout fibré holomorphe vectoriel sur $X$ est algébrique localement
trival.

\LayerTarget{tohoku:remark:3.4.2:3:p1}
3. La situation $u:P\longrightarrow A$ ($u$, hommomorphisme croisé de
faisceaux) décrite plus haut se rencontre dans diverses situations
intéressantes, par exemple : $X$ est une variété holomorphe, $G$ un groupe de
Lie complexe ayant l'algèbre de Lie $V$, $P$ est le faisceau des germes
d'application holomorphes de $X$ dans $G$, $A$ le faisceau des germes de
1-formes différentielles holomorphes sur $X$ à valeurs dans $V$, sur lequel
$P$ opère à l'aide de la représentation adjointe, et on pose

\LayerTarget{tohoku:eq:adjoint-logarithmic-derivative}
\[
u(g)=(dg)g^{-1}.
\]
Ceci permet, à tout espace fibré holomorphe $E$ sur $X$ de groupe structural
$G$, de faire correspondre

\LayerTarget{tohoku:def:adjoint-obstruction-class}
une classe
\[
d(E)\in H^1\bigl(X,\Omega^1(\operatorname{ad}(E))\bigr),
\]
où $\operatorname{ad}(E)$ désigne le fibré vectoriel (de fibre $V$)
\og adjoint\fg{} de $E$ et $\Omega^1(\operatorname{ad}(E))$ le faisceau des
germes de 1-formes différentielles holomorphes à valeurs dans
$\operatorname{ad}(E)$.

\LayerTarget{tohoku:construction:constant-transition-reduction:p1}
Le noyau de $u$ étant le sous-faisceau de $P$ formé des germes d'applications
constantes de $X$ dans $G$, on voit que l'annulation de $d(E)$ est une
condition nécessaire pour que le

\TohokuContinuation{166}
